% Deterministically converted from the hash-pinned frozen manuscript.
% Responsible human author: Carptopus.
\documentclass[11pt]{article}
\usepackage[a4paper,margin=25mm]{geometry}
\usepackage[T1]{fontenc}
\usepackage[utf8]{inputenc}
\usepackage{lmodern,microtype}
\usepackage{amsmath,amssymb,amsthm,mathtools}
\usepackage{needspace,xcolor,xurl,fancyhdr}
\usepackage[colorlinks=true,linkcolor=blue!55!black,citecolor=blue!55!black,urlcolor=blue!65!black]{hyperref}

\newtheorem{theorem}{Theorem}[section]
\newtheorem{proposition}[theorem]{Proposition}
\newtheorem{lemma}[theorem]{Lemma}

\pagestyle{fancy}
\fancyhf{}
\fancyhead[L]{Carptopus}
\fancyhead[R]{Effective eventual monotonicity}
\fancyfoot[C]{\thepage}
\setlength{\headheight}{14pt}
\setlength{\parindent}{0pt}
\setlength{\parskip}{0.42em}
\setlength{\emergencystretch}{4em}
\urlstyle{same}
\allowdisplaybreaks[1]
\title{Effective eventual monotonicity on the $(2,2,m)$ tensor slice}
\author{Carptopus\\\href{mailto:carptopus@163.com}{carptopus@163.com}}
\date{14 October 2026}
\hypersetup{pdfauthor={Carptopus},pdftitle={Effective eventual monotonicity on the (2,2,m) tensor slice}}
\begin{document}
\maketitle
\begin{center}
\fcolorbox{black!35}{black!4}{\parbox{0.86\textwidth}{\centering\small Internal proof audit passed. External peer review has not been completed.}}
\end{center}

\begin{abstract}

Let $a(n_1,\ldots,n_p)$ denote the expected number of real critical rank-one approximations to a standard Gaussian tensor of format $n_1\times\cdots\times n_p$, in the normalization of Draisma and Horobe\c{t}. They conjectured that, beyond the boundary format, this expectation does not increase when the largest format parameter is increased.  We prove an effective eventual form on the first open infinite slice:

\[
a(2,2,m)<a(2,2,m-1)
\qquad\text{for every integer }m\ge 3\times10^{103}+3.
\]

The proof first converts the Kac--Rice integral into an absolute determinant of a normalized random walk in $2\times2$ symmetric matrices.  Its exact three-dimensional moment-generating function permits a weighted Edgeworth expansion.  We certify positivity of the first even correction by symbolic reduction and outward-rounded interval integration.  Explicit low-, middle-, and high-frequency estimates then give a weighted Fourier remainder of order $r^{-5/2}$ with a numerical constant, yielding the displayed threshold.

This is an effective asymptotic tail theorem.  It does not settle the finite range below the threshold, including the smallest open comparison, and hence does not resolve the full monotonicity conjecture.

\end{abstract}

\section{Introduction}

Draisma and Horobe\c{t} studied the number of real critical points of the distance from a general tensor to the variety of rank-one tensors \hyperlink{ref1}{[1]}.  Because the real count depends on the tensor, they averaged it over the standard Gaussian distribution.  Their Conjecture 1.3 states that, if

\[
n_p-1>\sum_{i=1}^{p-1}(n_i-1),
\]

then replacing $n_p$ by $n_p-1$ does not decrease the expectation.  The matrix and several small-format cases are known, while the comparison between the formats $(2,2,3)$ and $(2,2,4)$ is the smallest case left open in their paper.

This paper proves monotonicity on an explicit tail of the entire $(2,2,m)$ slice.  The threshold is intentionally conservative: its role is to make the asymptotic argument completely effective, not to approximate the first true index of monotonicity.

In the sources checked through 14 October 2026, including the closest complex count and determinant-asymptotic results \hyperlink{ref3}{[3,4]} and the current problem and preprint indexes \hyperlink{ref5}{[5,6]}, we found no prior proof of this effective tail statement.  This is a bounded literature finding, not a claim of exhaustive global priority.  The random-matrix reduction below is a specialization and reorganization of the general Kac--Rice formulas in \hyperlink{ref1}{[1,2]}; the Edgeworth, Fourier, Plancherel, and Cauchy tools are standard.

\Needspace{5\baselineskip}
\begin{theorem}
For every integer

\[
m\ge 3\times10^{103}+3,
\]

one has

\[
a(2,2,m)<a(2,2,m-1).
\]

\end{theorem}

The theorem leaves every smaller parameter open unless it was already covered by prior work.  In particular, it says nothing about the comparison $a(2,2,4)\le a(2,2,3)$.

The proof has four parts.  Section 2 derives one common random-matrix model for the whole slice.  Sections 3 and 4 identify an exact transform and the positive first correction in its Edgeworth expansion.  Section 5 makes the Fourier remainder effective in three frequency regimes.  Section 6 combines the constants.  Section 7 records the computer-assisted certificates and their scope.

\section{A random walk in symmetric matrices}

Let

\[
J=\begin{pmatrix}0&1\\1&0\end{pmatrix},
\qquad P=v_0v_0^{\mathsf T},
\]

where $v_0\sim N(0,I_2)$.  Let $(v_i,y_i,z_i)_{i\ge1}$ be independent, with $v_i\sim N(0,I_2)$ and $y_i,z_i\sim N(0,1)$, and put

\[
K_i=v_iv_i^{\mathsf T}-y_i^2I_2+y_iz_iJ.                 \tag{2.1}
\]

For $r\ge1$ define

\[
I_r=\mathbb E\left|\det\left(P+\sum_{i=1}^rK_i\right)\right|. \tag{2.2}
\]

\Needspace{5\baselineskip}
\begin{proposition}
For every integer $m\ge3$,

\[
a(2,2,m)=\frac{\pi}{2(m-2)}I_{m-2}.                     \tag{2.3}
\]

\end{proposition}

\begin{proof}
Specialize \hyperlink{ref1}{[1, Theorem 1.1]} to block sizes $1,1,m-1$.  If $t,s$ are independent standard normal scalars and $x,y$ are independent standard normal vectors in $\mathbb R^{m-1}$, all mutually independent, the formula is

\[
a(2,2,m)=
\frac{(2\pi)^{3/2}}{2^{(m+4)/2}\Gamma(m/2)}
\mathbb E|\det C_m|,                                    \tag{2.4}
\]

where

\[
C_m=
\begin{pmatrix}
t&s&x^{\mathsf T}\\
s&t&y^{\mathsf T}\\
x&y&tI_{m-1}
\end{pmatrix}.                                           \tag{2.5}
\]

Put $r=m-2$.  Schur complementation, initially for $t\ne0$ and then by polynomial continuity, gives

\[
\det C_m=t^{r-1}
\det\!\left(t^2I_2+tsJ-GG^{\mathsf T}\right),           \tag{2.6}
\]

where $G$ is the $2\times(r+1)$ matrix with rows $x^{\mathsf T}$ and $y^{\mathsf T}$.  Tilting a standard normal scalar by $|t|^{r-1}$ makes $|t|$ chi distributed with $r$ degrees of freedom; its sign remains uniform and independent.  Hence, for independent standard normal vectors $u,z\in\mathbb R^r$,

\[
(t^2,ts)\quad\text{under the tilted law}
\ \overset d=\ (\|u\|^2,u\mathbin\cdot z),              \tag{2.7}
\]

where the sign of $t$ is absorbed into $s$.  The columns of $G$ are $r+1$ independent $N(0,I_2)$ vectors.  Negating the $2\times2$ matrix inside the determinant does not change its determinant.  Therefore (2.1)--(2.2) and (2.6)--(2.7) yield

\[
\mathbb E|\det C_m|
=\mathbb E|N(0,1)|^{r-1}I_r
=\frac{2^{(r-1)/2}\Gamma(r/2)}{\sqrt\pi}I_r.            \tag{2.8}
\]

Substituting $m=r+2$ and (2.8) into (2.4) gives $a(2,2,m)=\pi I_r/(2r)$, which is (2.3).

This is also the Wishart--Hessian decomposition obtained by specializing the two-dimensional chi-field Kac--Rice formula of \hyperlink{ref2}{[2]}.

\end{proof}

Consequently Theorem 1.1 is equivalent to the following statement.

\Needspace{5\baselineskip}
\begin{theorem}
For every integer $r\ge3\times10^{103}$,

\[
\frac{I_{r+1}}{r+1}<\frac{I_r}{r}.                       \tag{2.9}
\]

\end{theorem}

\section{Exact transform and whitening}

Write a symmetric matrix as $\bigl(\begin{smallmatrix}a&c\\c&b\end{smallmatrix}\bigr)$ and introduce the coordinates

\[
x_0=\frac{a+b}{2\sqrt3},\qquad
x_1=\frac{a-b}{2},\qquad
x_2=\frac c{\sqrt2}.                                    \tag{3.1}
\]

Let $X$ and $R$ be the vectors obtained from a single increment $K_i$ and the initial matrix $P$ in these coordinates.  For $s=(s_0,s_1,s_2)$ set

\[
\begin{aligned}
d_1(s)&=\left(1-\frac{s_0}{\sqrt3}\right)^2-s_1^2-\frac{s_2^2}{2},\\
d_2(s)&=1+\frac{2s_0}{\sqrt3}-\frac{s_2^2}{2}.
\end{aligned}                                            \tag{3.2}
\]

The Gaussian quadratic-form identity gives

\[
M_X(s)=\{d_1(s)d_2(s)\}^{-1/2},\qquad
M_R(s)=d_1(s)^{-1/2}.                                    \tag{3.3}
\]

Define

\[
Y_r=\frac{R+X_1+\cdots+X_r}{\sqrt r}.
\]

Then

\[
M_{Y_r}(s)=d_1(s/\sqrt r)^{-(r+1)/2}
d_2(s/\sqrt r)^{-r/2}.                                  \tag{3.4}
\]

Moreover, with

\[
q(x)=3x_0^2-x_1^2-2x_2^2,
\]

the coordinate change (3.1) gives

\[
\frac1r\det\left(P+\sum_{i=1}^rK_i\right)=q(Y_r),
\qquad
\frac{I_r}{r}=\mathbb E|q(Y_r)|.                        \tag{3.5}
\]

The order-zero term of the logarithm in (3.4) is $|s|^2/2$, so the increments are whitened: the limiting covariance is the identity.

\section{The positive first correction}

Let $\phi$ be the standard Gaussian density on $\mathbb R^3$.  Expanding the exact transform (3.4) through order $r^{-2}$ and translating monomials by the Fourier--Hermite correspondence gives polynomials $P_1,P_2,P_3,P_4$ such that formally

\[
p_r(x)=\phi(x)\left(1+\sum_{j=1}^4r^{-j/2}P_j(x)\right)
+\rho_r(x),                                             \tag{4.1}
\]

where $p_r$ is the density of $Y_r$.  The parity of $P_j$ is the parity of $j$.  Since $|q|\phi$ is even, the odd terms integrate to zero.  Thus

\[
\frac{I_r}{r}=L+\frac Cr+\frac D{r^2}+e_r,              \tag{4.2}
\]

where

\[
L=\int|q|\phi,\qquad
C=\int|q|\phi P_2,\qquad
D=\int|q|\phi P_4.                                     \tag{4.3}
\]

The first even polynomial is

\[
\begin{aligned}
P_2={}&x_0^6/54-x_0^4x_1^2/9-x_0^4/36+x_0^2x_1^4/6
+x_0^2x_1^2/2+7x_0^2x_2^2/12-5x_0^2/4\\
&+x_1^4/12+x_1^2x_2^2/4-17x_1^2/12
+x_2^4/8-4x_2^2/3+121/72.
\end{aligned}                                      \tag{4.4}
\]

\Needspace{5\baselineskip}
\begin{lemma}
The coefficient $C$ satisfies

\[
0.0853526992029<C<0.0864503864561.                       \tag{4.5}
\]

\end{lemma}

\begin{proof}
Integrate the Gaussian radial variable in (4.3) exactly.  On the unit sphere, write $x_0/|x|=u$ and set $t=\sin^2\theta$ in the remaining angular coordinate.  Since

\[
q(x)/|x|^2=(4+t)u^2-(1+t),
\]

the only zero for $u\in[0,1]$ is $u_*(t)=\sqrt{(1+t)/(4+t)}$.  Splitting the absolute value at $u_*(t)$ and integrating in $u$ symbolically yields

\[
C=\frac2\pi\int_0^{\pi/2}H(\sin^2\theta)\,d\theta,      \tag{4.6}
\]

where

\[
\begin{aligned}
H(t)=\frac1{8(t+4)^{7/2}}\bigl(&96t^6\sqrt{t+1}-48t^6\sqrt{t+4}
+1088t^5\sqrt{t+1}-472t^5\sqrt{t+4}\\
&+3360t^4\sqrt{t+1}-1098t^4\sqrt{t+4}
-1488t^3\sqrt{t+1}+1413t^3\sqrt{t+4}\\
&-2896t^2\sqrt{t+1}+4604t^2\sqrt{t+4}
+960t\sqrt{t+1}-2832t\sqrt{t+4}\\
&-192\sqrt{t+4}\bigr).
\end{aligned}                                            \tag{4.7}
\]

An outward-rounded interval Riemann sum over 20,000 equal subintervals of $[0,\pi/2]$ gives (4.5).  The symbolic calculation verifies that (4.7) is exactly the reduction of (4.4); the interval calculation supplies the strict sign.

\end{proof}

\section{An effective weighted Fourier remainder}

We now replace the formal remainder in (4.1) by a fully explicit bound.  All Fourier transforms use the convention for which inversion has factor $(2\pi)^{-3}$.

\subsection{Global modulus bounds}

For $u=(a,b,c)$ and $R=|u|$, direct multiplication gives

\[
\operatorname{Re}(d_1(iu)d_2(iu))
=1+R^2+\frac{b^2c^2}{2}+\frac{c^4}{4}-\frac{a^2c^2}{6}. \tag{5.1}
\]

It follows that

\[
|M_X(iu)|\le e^{-23R^2/94}\quad(R\le1),                 \tag{5.2}
\]

and

\[
|M_X(iu)|\le(1+R^2/2)^{-1/4}\quad(R\ge1).              \tag{5.3}
\]

For (5.2), use $a^2c^2\le R^4/4$ in (5.1).  For (5.3), minimize $|d_1(iu)|^2$ in $a^2$ and split according as $b^2\ge R^2/2$ or $a^2+c^2\ge R^2/2$.  In both cases the other determinant factor has modulus at least one.

\subsection{Low frequencies}

Put $\varepsilon=r^{-1/2}$.  If $|u|\le1/4$, the logarithm of the increment transform has a degree-six truncation $T_X$ satisfying

\[
|\log M_X(iu)-T_X(u)|\le4|u|^7.                         \tag{5.4}
\]

The initial factor has a degree-four truncation $T_R$ satisfying

\[
|\log M_R(iu)-T_R(u)|\le3|u|^5.                         \tag{5.5}
\]

Hence, for $|\xi|\le\sqrt r/4$,

\[
\left|\log M_R(i\xi/\sqrt r)+r\log M_X(i\xi/\sqrt r)
-T_R(\xi/\sqrt r)-rT_X(\xi/\sqrt r)\right|
\le r^{-5/2}(3|\xi|^5+4|\xi|^7).                      \tag{5.6}
\]

Exponentiating through order $\varepsilon^4$ and using (5.2) yields an explicit $r^{-5/2}$ Gaussian envelope.  A complex-polydisc Cauchy estimate with radius $1/[64(1+|\xi|)]$ upgrades it to all derivatives of order at most four.  For $r\ge2^{16}$,

\[
\left\|(1-\Delta)^2\widehat\rho_r\right\|_{L^2(
|\xi|\le r^{1/16})}
\le B_{\rm low}r^{-5/2},                                \tag{5.7}
\]

where

\[
B_{\rm low}<1.220\times10^{50}.                         \tag{5.8}
\]

One explicit ledger is

\[
\begin{aligned}
B_{\rm low}&=3.2\times10^{19}\sqrt J,\\
J&=4\pi2^{67}\left[\frac12(0.94)^{-3/2}\Gamma(3/2)
+\frac12(0.94)^{-71/2}\Gamma(71/2)\right].
\end{aligned}                                      \tag{5.9}
\]

\subsection{Middle and high frequencies}

Let $A=r^{1/16}$.  For $A<|\xi|\le\sqrt r$, a complex polydisc of radius $1/(64\sqrt r)$ remains in the zero-free region of both determinant factors. The logarithmic displacement is below $1/4$.  Cauchy's estimate and (5.2) therefore give

\[
|(1-\Delta)^2\widehat\rho_r(\xi)|
\le10^{30}r^6(1+|\xi|)^{40}e^{-23|\xi|^2/94}.           \tag{5.10}
\]

For $|\xi|=\sqrt r\,u$ with $u\ge1$, use a polydisc of radius $1/[4096\sqrt r(1+u)^2]$.  Its logarithmic displacement is below $1/512$, and (5.3) gives

\[
|(1-\Delta)^2\widehat\rho_r(\xi)|
\le10^{30}r^6(1+u)^{40}(1+u^2/2)^{-r/4}.                \tag{5.11}
\]

The gradient, zero-free, branch, and Cauchy-radius bounds leading to (5.10)--(5.11) are finite polynomial inequalities; they are included in the verification package.  Integrating (5.10) and (5.11) gives, for $r\ge2^{128}$,

\[
\left\|(1-\Delta)^2\widehat\rho_r\right\|_{L^2(
A<|\xi|\le\sqrt r)}\le r^{-5/2},                        \tag{5.12}
\]

and

\[
\left\|(1-\Delta)^2\widehat\rho_r\right\|_{L^2(
|\xi|\ge\sqrt r)}\le r^{-5/2}.                         \tag{5.13}
\]

For example, after squaring (5.11),

\[
\int_1^\infty u^2(1+u)^{80}(1+u^2/2)^{-r/2}\,du
\le(3/2)^{-(r/2-43)}\frac{2^{123}}3.                   \tag{5.14}
\]

The logarithmic derivative of the resulting high-frequency bound relative to $r^{-5/2}$ is

\[
\frac{37}{4}-\frac r4\log(3/2)<0                       \tag{5.15}
\]

throughout this range.

\subsection{From Fourier space to the weighted expectation}

Cauchy--Schwarz and Plancherel give

\[
\int_{\mathbb R^3}|q(x)|\,|\rho_r(x)|\,dx
\le\frac{3\sqrt5}{8\sqrt\pi}
\left\|(1-\Delta)^2\widehat\rho_r\right\|_2.          \tag{5.16}
\]

Here we used $|q(x)|\le3|x|^2$ and

\[
\int_{\mathbb R^3}\frac{|x|^4}{(1+|x|^2)^4}\,dx
=\frac{5\pi^2}{8}.
\]

Combining (5.7), (5.12), and (5.13), and rounding outward, gives

\[
|e_r|\le Kr^{-5/2},\qquad K<5.78\times10^{49}.          \tag{5.17}
\]

The exact second even Edgeworth polynomial has 51 monomials, all even and of total degree at most twelve.  Termwise Gaussian moment bounds give

\[
|D|\le\frac{540373}{144}<3753.                           \tag{5.18}
\]

\section{Proof of the main theorem}

From (4.2), (4.5), (5.17), and (5.18),

\[
\frac{I_r}{r}-\frac{I_{r+1}}{r+1}
\ge \frac{C}{2r^2}-\frac{3|D|}{r^3}-\frac{2K}{r^{5/2}}. \tag{6.1}
\]

The right-hand side is positive whenever both inequalities

\[
r>\frac{12|D|}{C},\qquad
\sqrt r>\frac{8K}{C}.                                  \tag{6.2}
\]

The strict numerical margins below verify these inequalities at the chosen integer threshold and therefore at every larger integer.

Outward-rounded interval arithmetic gives

\[
\frac{12|D|}{C}<5.28\times10^5,
\qquad
\left(\frac{8K}{C}\right)^2
<2.94\times10^{103}<3\times10^{103}.                   \tag{6.3}
\]

Therefore (2.9) holds for every integer $r\ge3\times10^{103}$.  Proposition 2.1 with $r=m-3$ translates this into Theorem 1.1.

\hfill$\square$

\section{Computer-assisted certificates}

The proof uses exact symbolic computation and outward-rounded interval arithmetic at three finite interfaces:

\begin{enumerate}
\item exact Gaussian quadratic-form determinants, whitening, and the two even Edgeworth polynomials;
\item symbolic reduction of $C$ to (4.6)--(4.7), followed by a 20,000-interval outward enclosure;
\item finite constant ledgers for the Cauchy bounds, Gaussian moments, and the threshold (6.3).
\end{enumerate}

The archived verification package is described in \texttt{verification/README.md}; \texttt{MANIFEST.sha256} fixes all nine checks, their shared exact Edgeworth generator, the pinned environment, and the analytic ledgers in \texttt{supplement/}.  The package also contains a destructive control showing that removing the initial rank-one factor changes the first correction, together with fixed-seed regression probes for the explicit analytic bounds.  Numerical quadrature used to calibrate $I_1$ and $I_2$ is excluded from the proof.

The interval implementation uses \texttt{mpmath.iv} endpoints without conversion to binary64 in the strict comparisons.  The proof does not infer a general statement from a finite parameter search.

\section{Scope and remaining problem}

The threshold in Theorem 1.1 is a safe effective bound produced by deliberately coarse global constants.  Improving those constants by a few orders of magnitude would not address the main remaining difficulty: the interval

\[
1\le r<3\times10^{103}
\]

is still far beyond direct enumeration.  A full solution of the $(2,2,m)$ slice will require a structural comparison, a propagation theorem, or a much stronger finite-core reduction.  No such statement is claimed here.

\section*{AI-assisted research and writing disclosure}
Carptopus is the sole author and is responsible for the mathematical claims, source attribution, revisions, and any public version of this manuscript. OpenAI Codex was used as an AI-assisted research and writing tool for literature organization, proof exploration, exact and interval verification, adversarial auditing, artifact generation, and manuscript drafting. Codex is not an author or responsible author.

\begin{thebibliography}{9}
\small\raggedright
\bibitem{ref1}\hypertarget{ref1}{}J. Draisma and E. Horobe\c{t}, \emph{The average number of critical rank-one approximations to a tensor}, Linear and Multilinear Algebra 64(12) (2016), 2498--2518. \url{https://doi.org/10.1080/03081087.2016.1164660}

\bibitem{ref2}\hypertarget{ref2}{}D. Marinucci and M. Stecconi, \emph{Critical points of chi-fields}, arXiv:2409.20129; ALEA 22 (2025), article 29. \url{https://doi.org/10.30757/ALEA.v22-29}

\bibitem{ref3}\hypertarget{ref3}{}S. Friedland and G. Ottaviani, \emph{The number of singular vector tuples and uniqueness of best rank-one approximation of tensors}, Foundations of Computational Mathematics 14 (2014), 1209--1242. \url{https://doi.org/10.1007/s10208-014-9194-z}

\bibitem{ref4}\hypertarget{ref4}{}T. Tao and V. Vu, \emph{A central limit theorem for the determinant of a Wigner matrix}, Advances in Mathematics 231 (2012), 74--101. \url{https://doi.org/10.1016/j.aim.2012.05.006}

\bibitem{ref5}\hypertarget{ref5}{}OpenProblemsInNLA, \emph{TR-16: Expected number of critical rank-one approximations of a Gaussian tensor}, status checked 14 October 2026. \url{https://github.com/ajt60gaibb/OpenProblemsInNLA/tree/main/tensor-computations/TR-16}

\bibitem{ref6}\hypertarget{ref6}{}OpenAI, \emph{math} public preprint index, commit \texttt{1581d4cbdd69a0faf34b2089614fe2a814316a6a}, checked 14 October 2026. \url{https://github.com/openai/math}

\end{thebibliography}
\end{document}
